\documentclass[11pt]{article}
\usepackage[T1]{fontenc}
\usepackage{lmodern,amsmath,amssymb,geometry,microtype,booktabs,longtable,array,xurl,hyperref,fancyvrb}
\geometry{letterpaper,margin=0.8in}
\hypersetup{colorlinks=true,linkcolor=blue,urlcolor=blue}
\setlength{\emergencystretch}{3em}
\providecommand{\tightlist}{\setlength{\itemsep}{0pt}\setlength{\parskip}{0pt}}
\fvset{fontsize=\small}
\title{A scaled dyadic Pell pair\\\large An explicit quartic in 212 unrestricted integers}
\author{Research continuation for ProveIt}
\date{26 October 2026}
\begin{document}
\maketitle
\begin{center}
\begin{tabular}{rrrrr}
Variables&Quadratic rows&Degree&Collected terms&Height\\\midrule
212&151&4&5,357&17,592,186,044,420
\end{tabular}
\end{center}
\section{A scaled dyadic Pell pair: 212 unrestricted integer
coordinates}\label{a-scaled-dyadic-pell-pair-212-unrestricted-integer-coordinates}

\subsection{Result and boundary}\label{result-and-boundary}

The literal equation \texttt{integer212/equation.json} is the sum of the
squares of 151 quadratic residuals in 212 independently unrestricted
integer variables. Its collected quartic has 5,357 nonzero terms and
height 17,592,186,044,420. The maximum residual coefficient is
4,194,304. The accompanying coefficient file lists every term. The
construction has an integer zero exactly when there is a finite
refutation in the fixed effective ZFC presentation. Consequently ZFC
proves its nonsolvability equivalent to Con(ZFC); this statement does
not assert either nonsolvability or consistency.

The immediate literal parent is source173's pre-optimization affine213
system, at ProveIt commit ecb9a4b4dda617512952525d50f17343700aac52. Its
complete unchanged original release is retained in
\texttt{inherited/Source173.zip}, including its ordinary proofs,
source-to-equation projections and independent verifier. The parent has
213 variables, 152 rows and, after its optional optimization, 5,402
quartic terms with height 2,417,851,639,242,452,488,945,685. Thus this
construction saves one coordinate, one row and 45 terms while lowering
that height by more than eleven decimal orders. The round22 215-variable
equation has much smaller coefficients and remains a useful incomparable
alternative. No global minimum or new proof-assistant compilation is
claimed.

The source173 anchor replacement, constructor-depth bound,
factorial-free compiler, Horn-flow argument and six exact projections
are credited, not new. The fourfold weak-frame pooling estimate comes
from round22; its proof is retained in
\texttt{credited/ROUND22\_AFFINE\_WEAK\_FRAME.md}. The new combination
is a scaled growth Pell pair, a weakened but sufficient cap on that
pair, and a nondegenerate affine growth base. Their unbounded proofs
follow here. Finite checks certify every literal algebraic change, not
these unbounded theorems. The inherited logical endpoints are
\texttt{KEndpoint.c222\_iff\_zf\_refutation} and
\texttt{ConstructibleChoice.zf\_refutation\_iff\_zfc\_refutation}, not a
machine execution or a specialization of Jones' universal polynomial.
The source closure is an ordinary-proof closure, not a complete
compilable Lean dependency tree.

\subsection{Changed semantic quantities and literal
coordinates}\label{changed-semantic-quantities-and-literal-coordinates}

Put D=2\^{}20=1048576. Keep the old coordinate t unrestricted, discard
the private coordinate that enforced t=16+4h\^{}2, and use

\begin{verbatim}
d=4t+2,          K=D t,          F=d^K.                 (1)
\end{verbatim}

The last notation is a theorem about a quadratic power certificate, not
an exponentiation operation in the literal equation. In particular no
sign of t or F is assumed here. The integer d is always even, nonzero
and nonunit.

Keep the thirteen principal Pell pairs with nominal indices

\begin{verbatim}
t, F+t, a,b,c,d_local,S,y+1,S',u,x,vq,C.
\end{verbatim}

Here d\_local is a local representative, distinct from growth base
d.~The reused anchor is the existing pair at C, and C is stored as twice
an integer. Replace the fourteenth, growth pair by three coordinates
w,z,Z representing

\begin{verbatim}
X=1+D w,       Y=D z,       Z=A z.
\end{verbatim}

Its two literal equations are

\begin{verbatim}
Z-Az=0,       2w+D(w^2+z^2-Z^2)=0.                    (2)
\end{verbatim}

Multiplying the second equation by D gives exactly
X\^{}2-(A\^{}2-1)Y\^{}2-1=0, modulo the first. Thus there is no
discarded divisibility condition or hidden extra variable. With
T=d\^{}2-2Ad+1, its power row is

\begin{verbatim}
F-1-Dw+DZ-dDz+T q_growth=0.                          (3)
\end{verbatim}

The sign of T is the opposite of the positive-form modulus 2Ad-d\^{}2-1
used below; divisibility is the same. All terms in (2)-(3) are at most
quadratic.

In the pooled orientation replace ONLY the growth-pair contribution

\begin{verbatim}
(X-1)^2+(Y-Dt)^2    by    w^2+(z-t)^2.              (4)
\end{verbatim}

The literal parent orientation has already had fourteen norm rows
subtracted. The producer removes its exact old growth contribution
2-2d\_old-2Z\_old+Z\_old\^{}2+D\^{}2t\^{}2-2DtY\_old, inserts the new
expression in (4), and later replays explicit integral row operations.
The checker reconstructs the raw orientation from all actual shifted
pairs before comparing this edit.

The anchor cap now uses z\^{}2, not Y\^{}2=D\^{}2z\^{}2, for this one
pair:

\begin{verbatim}
Y_C^2=C^2+1+sum_(other 12 nonanchor pairs) Y_i^2
             +z^2+r^2+s^2+v(v+1).                 (5)
\end{verbatim}

This is a semantic weakening, justified by the scaled-cap argument
below; it is not an asserted pointwise equivalence of arbitrary old
witness tuples. The literal cap row is unchanged because its former
ordinate coordinate now stores z.

\subsection{Weighted weak frame, parity and quantitative
growth}\label{weighted-weak-frame-parity-and-quantitative-growth}

All 35 actual source atoms (including duplicates) and all 48
base/index/output ports of the sixteen requests are explicitly listed in
the certificate. They include d=4t+2, K=Dt, C=2c0, F+t, q,q+1, BH, and
every power output. They are affine forms L=c+sum a\_j u\_j in
arithmetic coordinates only. Define w\_j as the largest
\textbar a\_j\textbar{} among these forms and set

\begin{verbatim}
W=22+sum_j w_j u_j^2.                              (6)
\end{verbatim}

In this ledger w\_t=D, w\_c0=2, and every other weight is 1;
\textbar c\textbar\textless=2. For every integer u,
\textbar u\textbar\textless=u\^{}2. Therefore each \textbar L\textbar{}
is at most 2+sum\_j w\_j u\_j\^{}2\textless W. This is an absolute
bound, NOT the false assertion L\^{}2\textless W. The complete atom and
coordinate census is checked independently.

Use the same fourfold Pell frame as round22:

\begin{verbatim}
U=4W(2v_frame-1),    V=(4W+1)U,    A^2+U^2-V^2=1.   (7)
\end{verbatim}

These are three quadratic rows using the existing coordinates. Since U,V
are even, A is odd. No sign of A is assumed. Classification of the Pell
solution (\textbar A\textbar,\textbar U\textbar) at center 4W+1,
together with U a nonzero multiple of 4W, gives an absolute index at
least 4W. Consequently

\begin{verbatim}
|A|>(4W)^(4W),   |A-1|>4W^3,   |A|>D+1.            (8)
\end{verbatim}

The last inequality already holds at W=22: 88\^{}4\textgreater D+1, a
fortiori 88\^{}88\textgreater D+1. For completeness choose the frame
index 4W. The familiar congruence Y\_(4W)(4W+1)=4W modulo 8W ensures
that U/(4W) is odd, supplying v\_frame. The frame is chosen only after
all finite arithmetic coordinates and W.

\subsection{Scaled cap and dyadic classification, before arithmetic
signs}\label{scaled-cap-and-dyadic-classification-before-arithmetic-signs}

The private anchor norm, squared-divisibility and comparison rows are
unchanged. Equation (5) gives
\textbar Y\_C\textbar\textgreater\textbar C\textbar{} and Y\_C!=0. The
source173 signed anchor proof therefore applies verbatim: classify its
absolute index i\textgreater0; the auxiliary comparison gives i=+/-C
modulo 2\textbar Y\_C\textbar; since
1\textless=i\textless=\textbar Y\_C\textbar{} and
\textbar C\textbar\textless\textbar Y\_C\textbar, C!=0 and
i=\textbar C\textbar. This uses no sign or interpreted arithmetic
meaning of C. In particular the possibly degenerate anchor index zero is
excluded before any comparison requiring a nontrivial pair.

Every unchanged pair has absolute index \textless\textbar C\textbar{} by
(5). For the scaled pair, (5) gives
\textbar z\textbar\textless\textbar Y\_C\textbar{} and hence
\textbar Y\textbar\textless D\textbar Y\_C\textbar. If its absolute
classified index were at least \textbar C\textbar+1, Pell recurrence at
a=\textbar A\textbar{} would give

\begin{verbatim}
|Y|>=Y_(|C|+1)(a)>a Y_|C|(a)>D|Y_C|,
\end{verbatim}

contradicting the preceding bound. Thus its absolute index is
\textless=\textbar C\textbar\textless W. Equality here is allowed;
strictness is unnecessary for orientation.

For odd a, every Pell abscissa X\_n(a) is odd. Reduction modulo 2 gives
Y\_n(a)=n modulo 2. Repeated doubling Y\_(2n)=2X\_nY\_n therefore
proves, for nonzero n,

\begin{verbatim}
v2(Y_n(a))=v2(n).                                 (9)
\end{verbatim}

Since D divides the scaled ordinate, the classified index is D*l for
some integer l. The zero ordinate has index zero and is included
separately. For D\textgreater=4, X\_(D*l)(a)=1 modulo D: use
X\_(2n)-1=2(a\^{}2-1)Y\_n\^{}2 and D/2 divides Y\_(D*l/2). Our physical
abscissa is 1 modulo D. It cannot be the negative of this canonical
abscissa, since +1 and -1 are distinct modulo D. At zero index, X=-1 is
excluded for the same reason.

These conclusions also handle A\textless0. The index D*l is even;
X\_(D*l)(A)=X\_(D*l)(\textbar A\textbar), and the ordinate sign chooses
signed l because Y\_(D*l)(A) changes sign with the sign of l. Thus the
physical pair is the canonical signed pair at the actual signed
parameter A.

Crucially we NEVER cancel D modulo the even integer A-1. Instead the
exact integer polynomial identities from repeated doubling give

Y\_(D*l)(A)/D=Y\_l(A) product\_(j=0,\ldots,19) X\_(2\^{}j*l)(A),
(X\_(D*l)(A)-1)/D = (D/2)(A\^{}2-1)(Y\_(D*l/2)(A)/(D/2))\^{}2.

The second identity is also valid at l=0. Evaluating these polynomials
at A=1 proves the genuine congruences

\begin{verbatim}
z=l mod(A-1),        w=0 mod(A-1).                 (10)
\end{verbatim}

They are valid at negative A and signed l as well.

The old orientation reductions for the other thirteen pairs have
abscissa residues +/-1 and signed ordinate residues of absolute value
\textless W. For the new pair (10) gives residues 0 and l, with
\textbar l\textbar\textless W. All fourteen ordinate errors have
absolute value \textless2W; all abscissa errors have absolute value at
most 2. The modified orientation thus has a nonnegative small residue E
with

\begin{verbatim}
0<=E<56W^2+56<4W^3<|A-1|       (W>=22).
\end{verbatim}

Divisibility by A-1 forces E=0. Every other pair is canonically oriented
at its nominal index. The new error gives l=t, so the scaled pair has
physical index Dt. Neither positivity of t nor any arithmetic power was
used.

\subsection{Power interpretation under the weak
frame}\label{power-interpretation-under-the-weak-frame}

For a pooled block of r\textless=15 requests at nonzero integer base n,
let its canonical pairs have indices k\_i, outputs m\_i, and Z\_i=A
Y\_i. The retained quadratic rows are

K\_n-2H-(2A-n)L\_n=0, S-nK\_n+L\_n=0, H=sum Y\_i(m\_i-X\_i),
S=sum((m\_i-X\_i+Z\_i)\^{}2-Y\_i\^{}2).

They imply that T\_n=2An-n\^{}2-1 divides
sum(m\_i-X\_i+Z\_i-nY\_i)\^{}2. Since T\_n=-1 modulo n, n is invertible
modulo T\_n.~Put b=W-1. Canonical signed-index Pell congruences,
multiplied by n\^{}(2b), give divisibility of

\begin{verbatim}
E=sum_i (m_i n^b-n^(b+k_i))^2.
\end{verbatim}

Every exponent is nonnegative and
\textbar n\textbar,\textbar m\_i\textbar,\textbar k\_i\textbar\textless=b.
For \textbar n\textbar\textgreater=2 each summand is
\textless=4b\^{}(4b); the same loose bound covers n=+/-1. Hence

\begin{verbatim}
0<=E<=60b^(4b)<(4W)^(4W)<|T_n|.
\end{verbatim}

The last inequality follows from
\textbar T\_n\textbar\textgreater=2\textbar A\textbar-(W-1)\^{}2-1\textgreater\textbar A\textbar{}
and (8). Thus E=0. At a nonunit base this forces k\_i\textgreater=0 and
m\_i=n\^{}k\_i. At a unit base it proves the signed unit power without
claiming a nonnegative index. The singleton requests have the same
conclusion, either directly or by squaring their divisibility error
(r=1). The scaled pair in (3) is just its physical canonical pair, so it
is covered. Conversely genuine powers give the two pooled quotient
witnesses by the identity S-2nH=T\_n L\_n and K\_n=2H+(2A-n)L\_n.~This
is an integral fresh witness construction, not an inverse on arbitrary
preliminary tuples.

\subsection{Noncircular horizon and signed arithmetic
bootstrap}\label{noncircular-horizon-and-signed-arithmetic-bootstrap}

The growth base d=4t+2 is always even and nonunit. The singleton lemma
first gives Dt\textgreater=0, hence t\textgreater=0, and F=d\^{}(Dt). If
t=0, then F=1 and y=F(2v+1) is odd. The endpoint equations reduce to

\begin{verbatim}
x=2+(y+1)u0,      x=1+(y+1)ut,
\end{verbatim}

which would make the even integer y+1 divide 1. This excludes t=0,
including y=-1. Therefore t\textgreater=1, d\textgreater=6 and
F\textgreater=6\^{}D\textgreater2\^{}32; F is positive and even.

Now follow the source173 sign recovery in its proper order. The raw row
B*eta=C+F+1 has odd right side, because C is even. Thus B!=0. The raw
radix multiplier H=2+4(C+a+b+c+d\_local+S+S'+u+vq) is 2 modulo 4,
regardless of representative signs, so
\textbar BH\textbar\textgreater=2. Its singleton gives x\textgreater=0
and q=(BH)\^{}x. If x=0, the last endpoint would make 1+(t+1)y divide
-1, impossible for t\textgreater=1 and \textbar y\textbar\textgreater=F.
Hence x\textgreater=1 and \textbar q\textbar\textgreater=2.

Interpret the q pool next. Its y+1 index is nonnegative. Since y is a
nonzero multiple of positive F, this gives y\textgreater=F. The first
endpoint forces u0\textgreater=0, and the last forces ut\textgreater=1.
In particular

\begin{verbatim}
x>=F+1,        x>=2+(t+1)y.                        (11)
\end{verbatim}

Now \textbar q\textbar\textgreater=2\^{}x and x\textgreater3, so
\textbar q+1\textbar\textgreater=2. Only now interpret the q+1 pool. All
nine actual representatives and their numerator powers are nonnegative;
H\textgreater=2. Normalized signed-remainder recovery authenticates
B=binom(F+t,t)\textgreater0. Therefore q\textgreater0 and
N=B\^{}x\textgreater0, as in the parent. No earlier step assumes these
last signs.

The source173 remainder argument still applies: its first-digit
coefficient bound (F+t)\^{}t\textless=2\^{}(t(F-1)) and
\textbar q\textbar\textgreater=2\^{}((t+1)F+2) hold for all the present
values. The common remainder norm, all ten digit rows, and all retained
strict local ranges are unchanged. The corner case y=F is excluded after
local recovery by the parent's positive-accumulator argument, not
prematurely.

\subsection{Small horizons, local label bounds and
soundness}\label{small-horizons-local-label-bounds-and-soundness}

Two old thresholds are explicitly replaced rather than silently reused.
First, Lambda\_t\textless=4\^{}(t-1), while F\textgreater=6\^{}(D
t)\textgreater16\^{}t for every t\textgreater=1. Thus the all-primes
binomial selector applies with radius R=F, including t=1. Second put
R0=2+t+x+y+F. Since F\textgreater=t, y\textgreater=F and
x\textgreater=2+2y, R0\textless3x. If
2\^{}s\textless=x\textless2\^{}(s+1), then s\textgreater=32 and

\begin{verbatim}
2^23 log2(3x)<2^23(s+3)<2^s<=x.
\end{verbatim}

The strict middle inequality holds at s=32 and persists by induction.
Consequently the inherited literal local-degree bound is
L\_local\textless R0\^{}(2\^{}23)\textless2\^{}x\textless=p\^{}(e x),
supplying exact local recovery.

For the deleted derived-label range, use the pinned source173 depth-19
constructor ledger with comparison base d=4t+2. At t\textgreater=2 we
have d\textgreater=10, so its estimate that each constructor on
arguments below s is \textless8s\^{}2 gives, by the usual exponent
recurrence, every label \textless d\^{}(D-1). The smallest horizon t=1
requires a separate argument, not a silent reuse of the hypothesis
d\textgreater=8: all local parameter leaves are zero, all constant
leaves are at most one, and we start with s0=2. Iterating sr+1=8sr\^{}2
gives sr=16\^{}(2\^{}r)/8. For every r\textgreater=0 this is less than
6\^{}(2\^{}(r+1)-1), since 16\textless36 and 1/8\textless1/6. At depth
r\textless=19 the label is therefore \textless6\^{}(D-1)=d\^{}(D-1) as
well. In all cases choose L=d\^{}D. Then

\begin{verbatim}
L>=t,      F=d^(Dt)=L^t.
\end{verbatim}

The source173 sharp product-separation lemma therefore applies even at
the smallest horizon. For clarity: unequal factor counts contradict
F+1\textgreater L\^{}t; for equal counts m\textless=t, each positive
elementary symmetric coefficient lies in {[}1,L\^{}m{]} subset
{[}1,F{]}, so differences have magnitude \textless F. The least nonzero
coefficient cannot vanish when evaluated at F. Unique factorization
recovers label multisets. Zero labels contribute units and are omitted.

The unchanged local-row, boundary, range and positive-accumulator
arguments now produce a balanced finite single-conclusion Horn flow. The
source173 Horn closure proof extracts a refutation even in the presence
of cycles. Its literal C222 calibration then yields a finite ZF
refutation and, by the retained constructible-choice equivalence, a
finite ZFC refutation. This completes soundness without circular
arithmetic positivity.

\subsection{Completeness and its order of
choices}\label{completeness-and-its-order-of-choices}

Start with a finite C222 refutation, all its local parameters and
padding. Choose an integer t\textgreater=1 sufficiently large to exceed
those finite parameters and its length. Set d=4t+2, K=Dt, F=d\^{}K.
Since v2(d)=1,

\begin{verbatim}
v2(F)=Dt>v2(t!),   v2(F+i)=v2(i) (1<=i<=t).
\end{verbatim}

The inherited beta CRT construction gives arbitrarily large positive y=F
a\^{}2 with a odd and divisible by the odd part of t!. It satisfies the
retained odd-linear equation with v=(a\^{}2-1)/2, and can dominate all
required finite accumulators. Choose an arbitrarily large positive beta
code x in its CRT class before defining N=B\^{}x or choosing the
independent representatives. All beta moduli are odd. The binomial
B=binom(F+t,t) is odd by the displayed valuations, so N\^{}2 is odd.

Choose the eight ordinary CRT representatives in their required residue
classes. Choose the remaining representative C even and cofinally large
in its own class modulo N\^{}2, exceeding every other nominal index
including Dt. This is possible because N\^{}2 is odd. Only then define H
and q=(BH)\^{}x and all associated true powers. The actual frame (6) is
computed after these arithmetic choices. No earlier modulus depends on
the enlarged C or q.

Choose the positive canonical fourfold frame A and every ordinary
principal pair. For the growth pair define

\begin{verbatim}
z=Y_(Dt)(A)/D,  w=(X_(Dt)(A)-1)/D,  Z=A z.
\end{verbatim}

The doubling identities above make these integers and make w and z-t
multiples of A-1. Equations (2)-(4), including the orientation quotient,
therefore have integral witnesses. Because C\textgreater Dt and C
exceeds the other indices, the parent cap slack remains nonnegative even
with the smaller scaled term z\^{}2. Its universal ternary
representation supplies the guard r\^{}2+s\^{}2+v(v+1). Every private
comparison witness is supplied by the same source173 construction: C
even gives odd abscissa and even ordinate; take comparison index C Y\_C,
use its odd squared quotient, and choose the comparison center by CRT.
This does not require the old horizon equation. Finally true powers
supply all singleton and pooled quotients and all unchanged local/Horn
wires. Apply the recorded integral coordinate changes and delete the
unused horizon-root slot. This yields a literal zero.

\subsection{Exact algebraic maps and
reproducibility}\label{exact-algebraic-maps-and-reproducibility}

The semantic step has the stated scope of fresh zero-existence
constructions, not an old-tuple bijection. Its physical pair map is
exactly

old\_d=1+D(w-Z), old\_Y=D z, old\_Z=D Z.

The weakened cap and discarded old horizon prohibit claiming that this
map alone lifts every new tuple to every old equation.
\texttt{base213/equation.json} is an explicitly marked staging array:
row130 and coordinate184 are inactive placeholders, not an additional
horizon condition. After the semantic step, every optional optimization
is genuinely bijective: each coordinate image is x\_i -\textgreater{}
x\_i+L(other coordinates), inverted by negating L; each row operation
R\_i -\textgreater{} R\_i+k R\_j is inverted by -k, in reverse order.
The private old root and its row are then omitted, and all remaining
indices are renumbered. \texttt{coordinate\_map.json} gives every exact
image and retained coordinate/row. \texttt{semantic\_certificate.json}
lists every changed literal coefficient, physical pair map, actual atom
and weight. No active certificate contains an absolute workspace path.
Historical path strings inside the unchanged source archive are
provenance only and are never used to locate new inputs.

\texttt{check.py} independently reconstructs the source identities,
complete edited system, every forward/inverse optimization, every
collected quartic term, all metrics and readable residuals.
\texttt{verify.py} runs it normally or with Python -O, copies the whole
package to a fresh temporary location, rebuilds there, and compares
every generated JSON byte. The optional \texttt{-\/-parent} also runs
the retained source173 verifier. None of these finite checks decides
consistency, searches for a zero, or substitutes for the ordinary proof.

\clearpage
\section{Complete literal residual appendix}
All $x_0,\ldots,x_{211}$ range independently over $\mathbb Z$. The equation is
$\sum_{i=0}^{150} R_i^2=0$. The following list is generated from, and checked
against, the complete machine-readable coefficient artifact.
\begingroup\scriptsize
\allowdisplaybreaks
% Entire literal residual list; generated from integer212/equation.json.
\begin{align*} R_{0}={}&x_{0} -x_{0}x_{4} -x_{4} +x_{7}
\end{align*}
\begin{align*} R_{1}={}&-x_{2} -x_{4}x_{5} +x_{6}x_{7}
\end{align*}
\begin{align*} R_{2}={}&-1 +x_{1} -x_{6}x_{7}
\end{align*}
\begin{align*} R_{3}={}&-2 x_{2}x_{211} -x_{3}x_{18} +x_{4} +2 x_{16}
\end{align*}
\begin{align*} R_{4}={}&-4 -3 x_{8} -x_{8}^{2} -2 x_{8}x_{9} -4 x_{9}
\\ &-x_{9}^{2} +x_{19}
\end{align*}
\begin{align*} R_{5}={}&x_{8} -x_{8}^{2} -8 x_{8}x_{19} -4 x_{8}x_{20} -16 x_{19}^{2}
\\ &-16 x_{19}x_{20} -4 x_{20}^{2} +x_{21}
\end{align*}
\begin{align*} R_{6}={}&1 +2 x_{21} -4 x_{21}^{2} +x_{22}
\end{align*}
\begin{align*} R_{7}={}&-2 -x_{9}^{2} -2 x_{9}x_{10} -x_{10}^{2} +x_{23}
\end{align*}
\begin{align*} R_{8}={}&1 -x_{8}^{2} -2 x_{8}x_{10} -x_{10}^{2} +x_{24}
\end{align*}
\begin{align*} R_{9}={}&-29 -61 x_{8} -35 x_{8}^{2} -96 x_{8}x_{10} -24 x_{8}x_{19}
\\ &-4 x_{8}x_{23} -24 x_{8}x_{24} +16 x_{9} -2 x_{9}^{2} -8 x_{9}x_{23}
\\ &-96 x_{10} -64 x_{10}^{2} -32 x_{10}x_{19} -32 x_{10}x_{24} -28 x_{19}
\\ &-4 x_{19}^{2} -8 x_{19}x_{24} +8 x_{23} -4 x_{23}^{2} -24 x_{24}
\\ &-4 x_{24}^{2} +x_{25}
\end{align*}
\begin{align*} R_{10}={}&-4 +10 x_{8} -2 x_{8}^{2} -8 x_{8}x_{9} -8 x_{8}x_{23}
\\ &-4 x_{25}^{2} +x_{26} -2 x_{43}
\end{align*}
\begin{align*} R_{11}={}&-4 -16 x_{10}^{2} -16 x_{10}x_{23} -16 x_{10}x_{24} -4 x_{23}^{2}
\\ &-8 x_{23}x_{24} -4 x_{24}^{2} +x_{27}
\end{align*}
\begin{align*} R_{12}={}&-7 +24 x_{10} -16 x_{10}^{2} -16 x_{10}x_{19} -32 x_{10}x_{24}
\\ &-16 x_{10}x_{27} +10 x_{19} -4 x_{19}^{2} -16 x_{19}x_{24} -8 x_{19}x_{27}
\\ &+24 x_{24} -16 x_{24}^{2} -16 x_{24}x_{27} +12 x_{27} -4 x_{27}^{2}
\\ &+x_{28}
\end{align*}
\begin{align*} R_{13}={}&-114 -127 x_{8} -37 x_{8}^{2} -24 x_{8}x_{9} -112 x_{8}x_{10}
\\ &-4 x_{8}x_{23} -24 x_{8}x_{24} -67 x_{9} -37 x_{9}^{2} -112 x_{9}x_{10}
\\ &-24 x_{9}x_{23} -4 x_{9}x_{24} -224 x_{10} -128 x_{10}^{2} -32 x_{10}x_{23}
\\ &-32 x_{10}x_{24} -16 x_{23} -4 x_{23}^{2} -40 x_{24} -4 x_{24}^{2}
\\ &+x_{29}
\end{align*}
\begin{align*} R_{14}={}&-30 -14 x_{8} -2 x_{8}^{2} -24 x_{8}x_{9} -32 x_{8}x_{10}
\\ &-8 x_{8}x_{23} -96 x_{9} -72 x_{9}^{2} -192 x_{9}x_{10} -48 x_{9}x_{23}
\\ &-128 x_{10} -128 x_{10}^{2} -64 x_{10}x_{23} -32 x_{23} -8 x_{23}^{2}
\\ &-4 x_{29}^{2} +x_{30}
\end{align*}
\begin{align*} R_{15}={}&-74 -42 x_{8} +8 x_{8}x_{10} +16 x_{8}x_{23} -84 x_{9}
\\ &+16 x_{9}x_{10} +32 x_{9}x_{23} +29 x_{10} -3 x_{10}^{2} -4 x_{10}x_{19}
\\ &-8 x_{10}x_{23} +21 x_{19} -8 x_{19}x_{23} +40 x_{23} -4 x_{23}^{2}
\\ &+x_{31} -x_{59}
\end{align*}
\begin{align*} R_{16}={}&-131 -184 x_{8} -64 x_{8}^{2} -160 x_{8}x_{10} -64 x_{8}x_{23}
\\ &-32 x_{8}x_{24} -32 x_{8}x_{31} -226 x_{10} -100 x_{10}^{2} -80 x_{10}x_{23}
\\ &-40 x_{10}x_{24} -40 x_{10}x_{31} -88 x_{23} -16 x_{23}^{2} -16 x_{23}x_{24}
\\ &-16 x_{23}x_{31} -46 x_{24} -4 x_{24}^{2} -8 x_{24}x_{31} -44 x_{31}
\\ &-4 x_{31}^{2} +x_{32}
\end{align*}
\begin{align*} R_{17}={}&-4 x_{19}^{2} +x_{20}
\end{align*}
\begin{align*} R_{18}={}&x_{8} -x_{8}^{2} +8 x_{8}x_{19} -4 x_{8}x_{20} -16 x_{19}^{2}
\\ &+16 x_{19}x_{20} -4 x_{20}^{2} +x_{33}
\end{align*}
\begin{align*} R_{19}={}&1 +2 x_{33} -4 x_{33}^{2} +x_{34}
\end{align*}
\begin{align*} R_{20}={}&-x_{9}^{2} -4 x_{9}x_{34} -4 x_{34}^{2} +x_{35}
\end{align*}
\begin{align*} R_{21}={}&-44 -52 x_{8} -16 x_{8}^{2} -64 x_{8}x_{9} +32 x_{8}x_{19}
\\ &-104 x_{9} -64 x_{9}^{2} +64 x_{9}x_{19} +50 x_{19} -16 x_{19}^{2}
\\ &+x_{36}
\end{align*}
\begin{align*} R_{22}={}&8 x_{9} -62 x_{9}^{2} -92 x_{9}x_{10} -32 x_{9}x_{23} -32 x_{9}x_{36}
\\ &+6 x_{10} -34 x_{10}^{2} -24 x_{10}x_{23} -24 x_{10}x_{36} -4 x_{23}^{2}
\\ &-8 x_{23}x_{36} +4 x_{36} -4 x_{36}^{2} +x_{37}
\end{align*}
\begin{align*} R_{23}={}&-90 -152 x_{8} -64 x_{8}^{2} -96 x_{8}x_{10} -32 x_{8}x_{24}
\\ &-32 x_{8}x_{37} -114 x_{10} -36 x_{10}^{2} -24 x_{10}x_{24} -24 x_{10}x_{37}
\\ &-38 x_{24} -4 x_{24}^{2} -8 x_{24}x_{37} -36 x_{37} -4 x_{37}^{2}
\\ &+x_{38}
\end{align*}
\begin{align*} R_{24}={}&-x_{10}^{2} -2 x_{10}x_{11} -x_{11}^{2} +x_{39}
\end{align*}
\begin{align*} R_{25}={}&-7 +8 x_{8} -3 x_{8}^{2} -8 x_{8}x_{9} -8 x_{8}x_{23}
\\ &-8 x_{8}x_{24} -4 x_{8}x_{39} -8 x_{9}x_{24} -2 x_{10}^{2} -8 x_{10}x_{39}
\\ &-8 x_{23}x_{24} +6 x_{24} -4 x_{24}^{2} -4 x_{39}^{2} +x_{40}
\\ &-x_{43}
\end{align*}
\begin{align*} R_{26}={}&-16 x_{8}^{2} -36 x_{8}x_{10} -32 x_{8}x_{24} -40 x_{8}x_{39} -16 x_{8}x_{40}
\\ &-4 x_{10} -22 x_{10}^{2} -32 x_{10}x_{24} -48 x_{10}x_{39} -16 x_{10}x_{40}
\\ &-2 x_{24} -16 x_{24}^{2} -32 x_{24}x_{39} -16 x_{24}x_{40} -4 x_{39}
\\ &-24 x_{39}^{2} -16 x_{39}x_{40} -4 x_{40}^{2} +x_{41}
\end{align*}
\begin{align*} R_{27}={}&x_{10} -25 x_{10}^{2} -40 x_{10}x_{24} -20 x_{10}x_{27} -16 x_{24}^{2}
\\ &-16 x_{24}x_{27} -4 x_{27}^{2} +x_{42}
\end{align*}
\begin{align*} R_{28}={}&-1 +8 x_{9} -4 x_{9}^{2} -8 x_{9}x_{23} +8 x_{23}
\\ &-4 x_{23}^{2} +x_{43}
\end{align*}
\begin{align*} R_{29}={}&-1 -16 x_{9}^{2} -8 x_{9}x_{10} -32 x_{9}x_{23} -16 x_{9}x_{43}
\\ &-x_{10}^{2} -8 x_{10}x_{23} -4 x_{10}x_{43} -16 x_{23}^{2} -16 x_{23}x_{43}
\\ &-4 x_{43}^{2} +x_{44}
\end{align*}
\begin{align*} R_{30}={}&3 +x_{9} -2 x_{9}x_{10} -4 x_{9}x_{44} -3 x_{10}^{2}
\\ &-8 x_{10}x_{44} -x_{23} -4 x_{44}^{2} +x_{45}
\end{align*}
\begin{align*} R_{31}={}&-1 -4 x_{45} -4 x_{45}^{2} +x_{46}
\end{align*}
\begin{align*} R_{32}={}&-1 -x_{10}^{2} -8 x_{10}x_{45} -4 x_{10}x_{46} -16 x_{45}^{2}
\\ &-16 x_{45}x_{46} -4 x_{46}^{2} +x_{47}
\end{align*}
\begin{align*} R_{33}={}&2 +x_{9} -x_{9}^{2} -4 x_{9}x_{10} -4 x_{9}x_{47}
\\ &-4 x_{10}^{2} -8 x_{10}x_{47} -4 x_{47}^{2} +x_{48}
\end{align*}
\begin{align*} R_{34}={}&1 -4 x_{42}^{2} -8 x_{42}x_{48} -4 x_{48}^{2} +x_{49}
\end{align*}
\begin{align*} R_{35}={}&1 +x_{8} -x_{8}^{2} -8 x_{8}x_{42} -4 x_{8}x_{49}
\\ &-16 x_{42}^{2} -16 x_{42}x_{49} -4 x_{49}^{2} +x_{50}
\end{align*}
\begin{align*} R_{36}={}&-1 -4 x_{50} -4 x_{50}^{2} +x_{51}
\end{align*}
\begin{align*} R_{37}={}&-1 -x_{10}^{2} -8 x_{10}x_{50} -4 x_{10}x_{51} -16 x_{50}^{2}
\\ &-16 x_{50}x_{51} -4 x_{51}^{2} +x_{52}
\end{align*}
\begin{align*} R_{38}={}&2 +x_{8} -x_{8}^{2} -4 x_{8}x_{10} -4 x_{8}x_{52}
\\ &-4 x_{10}^{2} -8 x_{10}x_{52} -4 x_{52}^{2} +x_{53}
\end{align*}
\begin{align*} R_{39}={}&-10 -12 x_{53} -4 x_{53}^{2} +x_{54}
\end{align*}
\begin{align*} R_{40}={}&-19 +4 x_{9} -4 x_{9}^{2} -8 x_{9}x_{23} -16 x_{9}x_{53}
\\ &-8 x_{9}x_{54} +4 x_{23} -4 x_{23}^{2} -16 x_{23}x_{53} -8 x_{23}x_{54}
\\ &-16 x_{53} -24 x_{53}^{2} -16 x_{53}x_{54} +6 x_{54} -4 x_{54}^{2}
\\ &+x_{55}
\end{align*}
\begin{align*} R_{41}={}&-10 +x_{9} -25 x_{9}^{2} -40 x_{9}x_{23} -20 x_{9}x_{55}
\\ &-16 x_{23}^{2} -16 x_{23}x_{55} -12 x_{48} -4 x_{48}^{2} -4 x_{55}^{2}
\\ &+x_{56}
\end{align*}
\begin{align*} R_{42}={}&21 +24 x_{48} -8 x_{48}^{2} -16 x_{48}x_{56} -2 x_{56}
\\ &-4 x_{56}^{2} +x_{57}
\end{align*}
\begin{align*} R_{43}={}&-1 +4 x_{8} +8 x_{9} +x_{10} -x_{10}^{2}
\\ &-4 x_{10}x_{57} -2 x_{19} -4 x_{57}^{2} +x_{58} +12 x_{59}
\\ &-4 x_{59}^{2}
\end{align*}
\begin{align*} R_{44}={}&-22 -38 x_{8} -16 x_{8}^{2} -64 x_{8}x_{9} +8 x_{8}x_{10}
\\ &+16 x_{8}x_{19} -76 x_{9} -64 x_{9}^{2} +16 x_{9}x_{10} +32 x_{9}x_{19}
\\ &+11 x_{10} -x_{10}^{2} -4 x_{10}x_{19} +19 x_{19} -4 x_{19}^{2}
\\ &+x_{59}
\end{align*}
\begin{align*} R_{45}={}&-4 +8 x_{8} +16 x_{9} -4 x_{19} -4 x_{58}^{2}
\\ &+24 x_{59} -8 x_{59}^{2} +x_{60}
\end{align*}
\begin{align*} R_{46}={}&-1 -4 x_{60} -4 x_{60}^{2} +x_{61}
\end{align*}
\begin{align*} R_{47}={}&-x_{2}x_{13} +x_{62}
\end{align*}
\begin{align*} R_{48}={}&x_{12} -x_{13} -2 x_{22}x_{62} -x_{63}
\end{align*}
\begin{align*} R_{49}={}&x_{12} -x_{13} -2 x_{26}x_{62} -x_{64}
\end{align*}
\begin{align*} R_{50}={}&x_{12} -x_{13} -2 x_{28}x_{62} -2 x_{62} -x_{65}
\end{align*}
\begin{align*} R_{51}={}&x_{12} -x_{13} -2 x_{30}x_{62} -x_{66}
\end{align*}
\begin{align*} R_{52}={}&x_{12} -x_{13} -2 x_{32}x_{62} +2 x_{62} -x_{67}
\end{align*}
\begin{align*} R_{53}={}&x_{12} -x_{13} -4 x_{34}x_{62} -2 x_{35}x_{62} -x_{68}
\end{align*}
\begin{align*} R_{54}={}&x_{12} -x_{13} -2 x_{38}x_{62} -x_{69}
\end{align*}
\begin{align*} R_{55}={}&2 x_{28}x_{62} -2 x_{41}x_{62} -2 x_{62} +x_{65} -x_{70}
\end{align*}
\begin{align*} R_{56}={}&-x_{8}x_{9} +x_{8}x_{10} +x_{9}x_{10} -x_{10}^{2} +x_{71}
\end{align*}
\begin{align*} R_{57}={}&1 +2 x_{11} +x_{11}^{2} -x_{71}^{2} +x_{72}
\end{align*}
\begin{align*} R_{58}={}&-x_{12}^{2} +2 x_{12}x_{13} +2 x_{12}x_{74} -x_{13}^{2} -2 x_{13}x_{74}
\\ &-x_{72}^{2} +x_{73} -x_{74}^{2}
\end{align*}
\begin{align*} R_{59}={}&-4 x_{60}x_{62} -2 x_{61}x_{62} +x_{74}
\end{align*}
\begin{align*} R_{60}={}&-x_{9}x_{62} -x_{13} +x_{75}x_{76} +x_{77}
\end{align*}
\begin{align*} R_{61}={}&-1 +8 x_{2} +4 x_{2}x_{8} +8 x_{2}x_{9} -2 x_{2}x_{19}
\\ &+x_{75}
\end{align*}
\begin{align*} R_{62}={}&-x_{8}x_{206} -x_{12} +x_{76}
\end{align*}
\begin{align*} R_{63}={}&x_{9}x_{206} -2 x_{19}x_{62} +2 x_{32}x_{62} +x_{67} +x_{78}
\end{align*}
\begin{align*} R_{64}={}&-x_{12}x_{63} +x_{13}x_{63} +x_{79}
\end{align*}
\begin{align*} R_{65}={}&-x_{64}x_{79} +x_{80}
\end{align*}
\begin{align*} R_{66}={}&-x_{65}x_{80} +x_{81}
\end{align*}
\begin{align*} R_{67}={}&-x_{66}x_{81} +x_{82}
\end{align*}
\begin{align*} R_{68}={}&-x_{67}x_{82} +x_{83}
\end{align*}
\begin{align*} R_{69}={}&-x_{68}x_{83} +x_{84}
\end{align*}
\begin{align*} R_{70}={}&-x_{69}x_{84} +x_{85}
\end{align*}
\begin{align*} R_{71}={}&-x_{70}x_{85} +x_{86}
\end{align*}
\begin{align*} R_{72}={}&-x_{77}x_{86} +x_{87}
\end{align*}
\begin{align*} R_{73}={}&-x_{78}x_{87} +x_{88}
\end{align*}
\begin{align*} R_{74}={}&-x_{88}^{2} +x_{89}
\end{align*}
\begin{align*} R_{75}={}&-x_{4} -2 x_{4}x_{16} +2 x_{16} +x_{90}
\end{align*}
\begin{align*} R_{76}={}&-x_{14}x_{90} +x_{91}
\end{align*}
\begin{align*} R_{77}={}&-2 x_{1}^{2} +2 x_{1}x_{12} +2 x_{1}x_{13} +2 x_{1}x_{91} +2 x_{1}x_{92}
\\ &-x_{12}^{2} -2 x_{12}x_{91} -x_{13}^{2} -2 x_{13}x_{92} +x_{17}x_{93}
\\ &-x_{73}x_{89} -x_{91}^{2} -x_{92}^{2}
\end{align*}
\begin{align*} R_{78}={}&-x_{4}x_{15} +x_{15} -x_{15}x_{90} +x_{92}
\end{align*}
\begin{align*} R_{79}={}&-x_{94}x_{97} +x_{97} +x_{99}
\end{align*}
\begin{align*} R_{80}={}&-x_{3}x_{120} +x_{95}x_{198} -x_{98} -x_{99}x_{120}
\end{align*}
\begin{align*} R_{81}={}&x_{100}x_{198} -x_{101}x_{123} -x_{102}
\end{align*}
\begin{align*} R_{82}={}&x_{103}x_{198} -x_{104}x_{123} -x_{105}
\end{align*}
\begin{align*} R_{83}={}&x_{106}x_{198} -x_{107}x_{123} -x_{108}
\end{align*}
\begin{align*} R_{84}={}&x_{109}x_{198} -x_{110}x_{123} -x_{111}
\end{align*}
\begin{align*} R_{85}={}&x_{112}x_{199} -x_{114}x_{123} -x_{115}
\end{align*}
\begin{align*} R_{86}={}&x_{116}x_{199} -x_{117}x_{123} -x_{118}
\end{align*}
\begin{align*} R_{87}={}&-x_{17}x_{120} +x_{123}
\end{align*}
\begin{align*} R_{88}={}&x_{119} -x_{121}x_{123} -x_{122}
\end{align*}
\begin{align*} R_{89}={}&-x_{123}x_{125} +x_{124} -x_{126}
\end{align*}
\begin{align*} R_{90}={}&-x_{123}x_{128} +x_{127}x_{198} -x_{129}
\end{align*}
\begin{align*} R_{91}={}&4 x_{130} -8 x_{130}x_{132} +x_{133}
\end{align*}
\begin{align*} R_{92}={}&-4 x_{130}x_{133} -x_{133} +x_{134}
\end{align*}
\begin{align*} R_{93}={}&-1 +x_{131}^{2} +x_{133}^{2} -x_{134}^{2}
\end{align*}
\begin{align*} R_{94}={}&1 -x_{131}x_{136} -x_{135} +x_{137}
\end{align*}
\begin{align*} R_{95}={}&-1 +2 x_{135} -x_{135}^{2} +2 x_{135}x_{137} +x_{136}^{2}
\end{align*}
\begin{align*} R_{96}={}&5 +16 x_{0} +16 x_{0}^{2} -8 x_{0}x_{131} -4 x_{131}
\\ &+x_{138}
\end{align*}
\begin{align*} R_{97}={}&-1 -4194304 x_{0}x_{181} +x_{2} +x_{138}x_{139} -1048576 x_{180}
\\ &-2097152 x_{181} +1048576 x_{182}
\end{align*}
\begin{align*} R_{98}={}&1 +x_{3}^{2} -2 x_{3}x_{131} +x_{140}
\end{align*}
\begin{align*} R_{99}={}&-x_{3}x_{172} +x_{17} +x_{140}x_{141} -x_{171}
\end{align*}
\begin{align*} R_{100}={}&-1 +2 x_{131}x_{183} -x_{142} -x_{183}^{2}
\end{align*}
\begin{align*} R_{101}={}&-1 +x_{94} +x_{142}x_{143} -x_{171} -x_{172}x_{183}
\end{align*}
\begin{align*} R_{102}={}&-x_{131}x_{145} -x_{144} +x_{146}
\end{align*}
\begin{align*} R_{103}={}&-1 -x_{144}^{2} +2 x_{144}x_{146} +x_{145}^{2}
\end{align*}
\begin{align*} R_{104}={}&-x_{131}x_{148} -x_{147} +x_{149}
\end{align*}
\begin{align*} R_{105}={}&-1 -x_{147}^{2} +2 x_{147}x_{149} +x_{148}^{2}
\end{align*}
\begin{align*} R_{106}={}&-x_{131}x_{151} -x_{150} +x_{152}
\end{align*}
\begin{align*} R_{107}={}&-1 -x_{150}^{2} +2 x_{150}x_{152} +x_{151}^{2}
\end{align*}
\begin{align*} R_{108}={}&-x_{131}x_{154} -x_{153} +x_{155}
\end{align*}
\begin{align*} R_{109}={}&-1 -x_{153}^{2} +2 x_{153}x_{155} +x_{154}^{2}
\end{align*}
\begin{align*} R_{110}={}&-x_{131}x_{157} -x_{156} +x_{158}
\end{align*}
\begin{align*} R_{111}={}&-1 -x_{156}^{2} +2 x_{156}x_{158} +x_{157}^{2}
\end{align*}
\begin{align*} R_{112}={}&-x_{131}x_{160} -x_{159} +x_{161}
\end{align*}
\begin{align*} R_{113}={}&-1 -x_{159}^{2} +2 x_{159}x_{161} +x_{160}^{2}
\end{align*}
\begin{align*} R_{114}={}&1 -x_{131}x_{163} -x_{162} +x_{164}
\end{align*}
\begin{align*} R_{115}={}&-1 +2 x_{162} -x_{162}^{2} +2 x_{162}x_{164} +x_{163}^{2}
\end{align*}
\begin{align*} R_{116}={}&-x_{131}x_{166} -x_{165} +x_{167}
\end{align*}
\begin{align*} R_{117}={}&-1 -x_{165}^{2} +2 x_{165}x_{167} +x_{166}^{2}
\end{align*}
\begin{align*} R_{118}={}&-x_{131}x_{169} -x_{168} +x_{170}
\end{align*}
\begin{align*} R_{119}={}&-1 -x_{168}^{2} +2 x_{168}x_{170} +x_{169}^{2}
\end{align*}
\begin{align*} R_{120}={}&1 -x_{131}x_{172} -x_{171} +x_{173}
\end{align*}
\begin{align*} R_{121}={}&-1 +2 x_{171} -x_{171}^{2} +2 x_{171}x_{173} +x_{172}^{2}
\end{align*}
\begin{align*} R_{122}={}&-x_{131}x_{175} -x_{174} +x_{176}
\end{align*}
\begin{align*} R_{123}={}&-1 -x_{174}^{2} +2 x_{174}x_{176} +x_{175}^{2}
\end{align*}
\begin{align*} R_{124}={}&-x_{131}x_{178} -x_{177} +x_{179}
\end{align*}
\begin{align*} R_{125}={}&-1 -x_{177}^{2} +2 x_{177}x_{179} +x_{178}^{2}
\end{align*}
\begin{align*} R_{126}={}&-x_{131}x_{181} +x_{182}
\end{align*}
\begin{align*} R_{127}={}&2 x_{180} +1048576 x_{180}^{2} +1048576 x_{181}^{2} -1048576 x_{182}^{2}
\end{align*}
\begin{align*} R_{128}={}&-23 -1048576 x_{0}^{2} -x_{1}^{2} -x_{2}^{2} -x_{3}^{2}
\\ &-x_{4}^{2} -2 x_{8} -x_{8}^{2} -x_{9}^{2} -x_{10}^{2}
\\ &-x_{11}^{2} -x_{12}^{2} -x_{13}^{2} -x_{14}^{2} -x_{15}^{2}
\\ &-2 x_{16}^{2} -x_{17}^{2} -x_{94}^{2} -x_{95}^{2} -x_{96}^{2}
\\ &-x_{100}^{2} -x_{103}^{2} -x_{106}^{2} -x_{109}^{2} -x_{112}^{2}
\\ &-x_{113}^{2} -x_{116}^{2} -x_{119}^{2} -x_{120}^{2} -x_{124}^{2}
\\ &-x_{127}^{2} +x_{130} -x_{183}^{2}
\end{align*}
\begin{align*} R_{129}={}&-2 x_{3} -4 x_{3}x_{8} -4 x_{3}x_{9} -4 x_{3}x_{10} -4 x_{3}x_{11}
\\ &-4 x_{3}x_{12} -4 x_{3}x_{13} -4 x_{3}x_{14} -4 x_{3}x_{15} -8 x_{3}x_{16}
\\ &+x_{183}
\end{align*}
\begin{align*} R_{130}={}&-x_{131}x_{185} +x_{194}
\end{align*}
\begin{align*} R_{131}={}&-1 +x_{184}^{2} +x_{185}^{2} -x_{194}^{2}
\end{align*}
\begin{align*} R_{132}={}&-x_{187}x_{188} +x_{195}
\end{align*}
\begin{align*} R_{133}={}&-1 +x_{186}^{2} +x_{187}^{2} -x_{195}^{2}
\end{align*}
\begin{align*} R_{134}={}&-x_{178}^{2} +x_{196}
\end{align*}
\begin{align*} R_{135}={}&x_{185} -2 x_{190}x_{196} -x_{196}
\end{align*}
\begin{align*} R_{136}={}&-1 -4 x_{185} -8 x_{185}x_{189} +x_{188}
\end{align*}
\begin{align*} R_{137}={}&-x_{131} -x_{184}x_{191} +x_{188}
\end{align*}
\begin{align*} R_{138}={}&-x_{179} -x_{184}x_{192} +x_{186}
\end{align*}
\begin{align*} R_{139}={}&-2 x_{16} -4 x_{178}x_{193} +x_{187}
\end{align*}
\begin{align*} R_{140}={}&9 +3 x_{0}^{2} +2 x_{0}x_{2} -2 x_{0}x_{136} -2 x_{0}x_{145}
\\ &-2 x_{0}x_{181} +x_{1}^{2} -2 x_{1}x_{172} +x_{2}^{2} -2 x_{2}x_{145}
\\ &+x_{4}^{2} -2 x_{4}x_{163} +x_{8}^{2} -2 x_{8}x_{148} +x_{9}^{2}
\\ &-2 x_{9}x_{151} +x_{10}^{2} -2 x_{10}x_{154} +x_{11}^{2} -2 x_{11}x_{157}
\\ &+x_{12}^{2} -2 x_{12}x_{160} +x_{13}^{2} -2 x_{13}x_{166} +x_{14}^{2}
\\ &-2 x_{14}x_{169} +x_{15}^{2} -2 x_{15}x_{175} -4 x_{16}x_{178} -x_{131}x_{197}
\\ &+x_{137}^{2} -2 x_{146} +x_{146}^{2} -2 x_{149} +x_{149}^{2}
\\ &-2 x_{152} +x_{152}^{2} -2 x_{155} +x_{155}^{2} -2 x_{158}
\\ &+x_{158}^{2} -2 x_{161} +x_{161}^{2} +x_{164}^{2} -2 x_{167}
\\ &+x_{167}^{2} -2 x_{170} +x_{170}^{2} +x_{173}^{2} -2 x_{176}
\\ &+x_{176}^{2} +2 x_{178}^{2} -2 x_{179} +x_{179}^{2} +x_{180}^{2}
\\ &+x_{197} -x_{200}^{2} -x_{201}^{2} -x_{202} -x_{202}^{2}
\end{align*}
\begin{align*} R_{141}={}&-x_{96}x_{198} +x_{120}
\end{align*}
\begin{align*} R_{142}={}&-x_{113}x_{199} +x_{120}
\end{align*}
\begin{align*} R_{143}={}&-1 -4 x_{16}^{2} -x_{136}^{2} -x_{145}^{2} -x_{148}^{2}
\\ &-x_{151}^{2} -x_{154}^{2} -x_{157}^{2} -x_{160}^{2} -x_{163}^{2}
\\ &-x_{166}^{2} -x_{169}^{2} -x_{172}^{2} -x_{175}^{2} +x_{178}^{2}
\\ &-x_{181}^{2} -x_{200}^{2} -x_{201}^{2} -x_{202} -x_{202}^{2}
\end{align*}
\begin{align*} R_{144}={}&-3 -4 x_{98}^{2} -4 x_{102}^{2} -4 x_{105}^{2} -4 x_{108}^{2}
\\ &-4 x_{111}^{2} -4 x_{115}^{2} -4 x_{118}^{2} +x_{120}^{2} -4 x_{122}^{2}
\\ &-4 x_{126}^{2} -4 x_{129}^{2} -x_{203}^{2} -x_{204}^{2} +3 x_{205}
\\ &-x_{205}^{2}
\end{align*}
\begin{align*} R_{145}={}&-x_{2}x_{12} +x_{206}
\end{align*}
\begin{align*} R_{146}={}&x_{94}x_{208} -2 x_{95}x_{145} -2 x_{100}x_{148} -2 x_{103}x_{151} -2 x_{106}x_{154}
\\ &-2 x_{109}x_{157} -2 x_{112}x_{160} -2 x_{116}x_{166} -2 x_{119}x_{169} -2 x_{124}x_{175}
\\ &-2 x_{127}x_{178} -2 x_{131}x_{208} +2 x_{145}x_{146} +2 x_{148}x_{149} +2 x_{151}x_{152}
\\ &+2 x_{154}x_{155} +2 x_{157}x_{158} +2 x_{160}x_{161} +2 x_{166}x_{167} +2 x_{169}x_{170}
\\ &+2 x_{175}x_{176} +2 x_{178}x_{179} +x_{207}
\end{align*}
\begin{align*} R_{147}={}&-10 -x_{94}x_{207} +x_{95}^{2} -2 x_{95}x_{144} +x_{100}^{2}
\\ &-2 x_{100}x_{147} +x_{103}^{2} -2 x_{103}x_{150} +x_{106}^{2} -2 x_{106}x_{153}
\\ &+x_{109}^{2} -2 x_{109}x_{156} +x_{112}^{2} -2 x_{112}x_{159} +x_{116}^{2}
\\ &-2 x_{116}x_{165} +x_{119}^{2} -2 x_{119}x_{168} +x_{124}^{2} -2 x_{124}x_{174}
\\ &+x_{127}^{2} -2 x_{127}x_{177} +2 x_{144}x_{146} +2 x_{147}x_{149} +2 x_{150}x_{152}
\\ &+2 x_{153}x_{155} +2 x_{156}x_{158} +2 x_{159}x_{161} +2 x_{165}x_{167} +2 x_{168}x_{170}
\\ &+2 x_{174}x_{176} +2 x_{177}x_{179} +x_{208}
\end{align*}
\begin{align*} R_{148}={}&x_{94}x_{210} -2 x_{96}x_{136} -2 x_{113}x_{163} -2 x_{120}x_{172} -2 x_{131}x_{210}
\\ &+2 x_{136} +2 x_{136}x_{137} +2 x_{163} +2 x_{163}x_{164} +2 x_{172}
\\ &+2 x_{172}x_{173} +x_{209} -x_{210}
\end{align*}
\begin{align*} R_{149}={}&-2 -x_{94}x_{209} +x_{96}^{2} -2 x_{96}x_{135} +x_{113}^{2}
\\ &-2 x_{113}x_{162} +x_{120}^{2} -2 x_{120}x_{171} +2 x_{135} +2 x_{135}x_{137}
\\ &+2 x_{162} +2 x_{162}x_{164} +x_{171}^{2} -x_{172}^{2} +x_{209}
\\ &+x_{210}
\end{align*}
\begin{align*} R_{150}={}&-1 -x_{2} -2 x_{2}x_{211} +x_{4}
\end{align*}

\endgroup
\end{document}
